\scanpage{30}
\[
\Sigma'(x)=\begin{cases}
R(x)&\text{si }x\in C(j)\cup G(j)\\
R(x)\cap\complement G(j)&\text{si }x\in\bar C(j)\cap\complement G(j)
\end{cases}
\]
(C'est une stratégie contre $j$, par [b) $\supset$]. Alors Pour $x\in C$, les conditions suivantes sont équivalentes

a) $x_0\notin G(j)$

b) $x_0$ est admissible contre $j$

c) $x_0$ est admissible pour $\Sigma'$ contre $j$.

\uline{Dém}\quad c) $\Longrightarrow$ b) trivial, b) $\Longrightarrow$ a) par le lemme, Prouvons a) $\Longrightarrow$ c). \unsure{Si on a} une partie $(x_0,x_1,\ldots)$ d'origine $x_0$ jouée en accord avec $\Sigma'$, alors par définition de $\Sigma$, on voit de proche en proche que $x_i\in\complement G(j)$ $\forall i$, donc la partie ne peut être gagnée par $j$ !

\uline{Corollaire}\quad Supposons $J=\{j,j'\}$, $j\ne j'$, et
\margin{$\big|$}
\[
G_0(j)\cap G_0(j')=\varnothing.
\]
Soit \deleted{$\unsure{N}$}
\[
N=\{x\in C\mid x\text{ admissible contre }j,\text{ et contre }j'\}.
\]
\margin{$\big|$}
Alors $C$ est réunion disjointe de $G(j)$, $G(j')$ et $N$.
% archive foliation: 29
